\chapter*{From source documents to transaction decisions}
\addcontentsline{toc}{chapter}{From source documents to transaction decisions}
\markboth{Process guide}{}
\label{map:overview}
\begingroup
\tikzset{
 mapbox/.style={draw=teal,fill=pale,rounded corners=2pt,align=center,
   text width=42mm,minimum height=16mm,inner sep=3mm,font=\small},
 mapwide/.style={mapbox,text width=108mm,minimum height=11mm,inner sep=2mm},
 mapresearch/.style={mapbox,draw=navy,dashed,fill=white},
 maparr/.style={-{Latex[length=2mm]},thick,draw=navy},
 mapnote/.style={font=\small,align=center,fill=white,inner sep=1mm}
}
\newcommand{\MapSubheading}[1]{\par\vspace{2mm}{\sffamily\bfseries\large #1\par}\vspace{1mm}}
\newcommand{\MapHeading}[2]{\phantomsection\label{#2}\MapSubheading{#1}}
\newenvironment{MapDiagram}{\par\noindent\begin{minipage}{\linewidth}\centering}{\par\end{minipage}\par}
\newcommand{\MapCaption}[1]{\par{\small\itshape #1\par}\vspace{2mm}}

A prospectus can permit a particular purchaser without establishing that a
bank may execute the purchase. The bank must also establish which laws apply,
who owns the account, whether sanctions restrict the transaction, and what its
own mandate and settlement arrangements permit. LegalMath keeps these questions
separate so that an answer to one cannot silently answer the others.

The running example compares two real contingent convertible securities, or
\emph{CoCos}: Barclays' February 2025 issue and Standard Chartered's January
2025 issue. These securities can absorb losses through conversion into shares
and other contractual or regulatory measures. The same hypothetical US
individual passes one selected purchaser condition and fails the other. Both
complete bank assessments remain qualified because essential evidence and
legal coverage are missing. The following charts trace that distinction.

There are two kinds of layer to keep apart. The \emph{software stages} turn
source documents and proposed meanings into executable calculations. The
\emph{seven assessment categories} organise the different legal and factual
questions about a transaction. Every category uses evidence and scoped rules;
it is not a separate programming language. The eight-step overview links to
both the worked CoCo case and the implementation details.

The first software stage, called the \emph{frontend}, proposes a meaning for the
passage and records its sources, assumptions and required facts. A
\emph{translator} then turns that fixed proposal into an executable language.
LegalMath has two such targets. \emph{RuleIR}, its own small rule language,
expresses conditions such as a financial threshold and the effect of missing
evidence. \emph{Catala} is a programming language designed for legal rules,
with a compiler that translates those rules into executable code. Both targets
can receive the same proposed meaning; Catala also supports the implemented
calculations over structured collections described in Chapter~\ref{ch:implementation}.

\emph{Stipula} answers a different question: how does a contract change when a
party takes an action? It describes parties, resources and states. For example,
a request available while an agreement is active may become unavailable after
withdrawal. The book studies Stipula's verification methods as a comparison;
LegalMath does not currently compile Stipula contracts.

\clearpage
\MapSubheading{Eight steps from documents to a qualified decision}
\begin{MapDiagram}
\begin{tikzpicture}[node distance=5mm]
\node[mapwide] (a) {\hyperref[map:source]{\textbf{1. Retain the source documents}}\\Prospectuses, laws, circulars and factual evidence};
\node[mapwide,below=of a] (b) {\hyperref[map:context]{\textbf{2. Fix context and applicable questions}}\\Entities, accounts, actions, dates and seven assessment categories};
\node[mapwide,below=of b] (c) {\hyperref[map:readings]{\textbf{3. Develop possible interpretations}}\\Duties, exceptions and required facts};
\node[mapwide,below=of c] (d) {\hyperref[map:challenge]{\textbf{4. Challenge the interpretations}}\\Omissions, alternative readings and revealing cases};
\node[mapwide,below=of d] (e) {\hyperref[map:rules]{\textbf{5. Construct the shared formal model}}\\Typed facts, scoped rules and explicit uncertainty};
\node[mapwide,below=of e] (f) {\hyperref[map:java]{\textbf{6. Generate executable targets and check behaviour}}\\RuleIR/Java and Catala execution tests};
\node[mapwide,below=of f] (g) {\hyperref[map:delivery]{\textbf{7. Assemble the transaction assessment}}\\All declared obligations, conditional results and remaining gaps};
\node[mapwide,below=of g] (h) {\hyperref[map:operation]{\textbf{8. Reassess when evidence or events change}}\\Invalidate stale findings; retain duties and earlier decisions};
\foreach \x/\y in {a/b,b/c,c/d,d/e,e/f,f/g,g/h}
  \draw[maparr] (\x)--(\y);
\draw[maparr] (d.east)--++(12mm,0)|- node[pos=.25,mapnote,rotate=90]{Revise} (c.east);
\draw[maparr] (f.west)--++(-12mm,0)|- node[pos=.25,mapnote,rotate=90]{Repair} (e.west);
\end{tikzpicture}
\end{MapDiagram}
\MapCaption{A failure or an unknown remains part of the assessment. Completing
the software steps does not establish complete legal coverage or permission to
execute a transaction.}

The next pages expand the \hyperref[map:layers]{seven assessment categories}
and follow the \hyperref[map:coco-terms]{two CoCo purchaser conditions} into the
\hyperref[map:coco-assessment]{complete bank assessment}. Later pages explain
each software step. Unsupported expressions prevent translation; missing facts
can prevent execution even when a program has compiled. A different proposed
interpretation remains a different model, with its own evidence.

\clearpage
\MapHeading{The seven layers of a transaction assessment}{map:layers}
Begin with a proposed action by a specified entity for a specified account.
The bank's brand does not identify its booking branch, capacity or the client's
status. Follow all seven categories: an adverse or unresolved finding remains
in the assessment while the other questions are considered.

\begin{MapDiagram}
\begin{tikzpicture}[node distance=4mm]
\node[mapwide,text width=132mm] (a) {\textbf{1. Scope: whose action, where and when?}\\Legal entities, branches, principal or intermediary capacity;\\jurisdictions, activity, assessment date and source editions};
\node[mapwide,text width=132mm,below=of a] (b) {\textbf{2. Sanctions: which restrictions reach the transaction?}\\Issuer and owner identity, designated programs, ownership,\\covered securities, licences and effective dates};
\node[mapwide,text width=132mm,below=of b] (c) {\textbf{3. Bank: may this entity provide this service?}\\Activity and licensing rules; selected US Bank Secrecy Act\\program controls where their applicability is established};
\node[mapwide,text width=132mm,below=of c] (d) {\textbf{4. Client: which account duties apply?}\\Beneficial owners, account purpose and required diligence;\\other anti-money-laundering and client requirements};
\node[mapwide,text width=132mm,below=of d] (e) {\textbf{5. Product: which terms and selling rules apply?}\\Prospectus, issuer law and loss absorption; purchaser restrictions;\\Hong Kong product and sophisticated-investor conditions};
\node[mapwide,text width=132mm,below=of e] (f) {\textbf{6. Internal policy: what does the actual mandate permit?}\\Versioned bank policies and account restrictions;\\missing private documents remain missing inputs};
\node[mapwide,text width=132mm,below=of f] (g) {\textbf{7. Operations: can this route satisfy its obligations?}\\Custody, intermediaries, currency and settlement;\\distinct rejection, blocking, reporting and disclosure duties};
\node[mapwide,text width=132mm,below=of g] (h) {\textbf{Retain every declared obligation and its evidence}\\Conditional result, conflict, missing input or unsupported rule};
\foreach \x/\y in {a/b,b/c,c/d,d/e,e/f,f/g,g/h}\draw[maparr](\x)--(\y);
\end{tikzpicture}
\end{MapDiagram}
\MapCaption{The current inventory has fourteen obligations across these seven
categories. Five have conditional rule models; nine retain missing rule families
or unresolved coverage. None may be removed by the caller.}

The five calculated obligations concern selected sanctions, bank program
controls, private-account diligence, sophisticated professional investor
conditions and reporting. Supporting calculations bring the collection to
eleven shared models. Complete discovery of applicable law remains unproved.
Section~\ref{sec:bank-compliance} explains the provisions and their limits.

\clearpage
\MapHeading{The CoCo example: from terms to purchaser conditions}{map:coco-terms}
Assume one individual, resident and located in the United States, buys
USD~200,000 principal for their own account and benefit. The individual is not
a \emph{qualified institutional buyer} (QIB), an institutional and account
category used by Rule~144A. Wealth alone does not turn an individual into a
QIB. Regulation~S supplies a different offshore distribution route; its
US-person and account definitions must also be considered.

\begin{MapDiagram}
\begin{tikzpicture}[node distance=8mm and 8mm]
\tikzset{pair/.style={mapbox,text width=62mm,minimum height=22mm}}
\node[pair] (a) {\textbf{Barclays PLC, 7.625\%}\\Supplement: 18 February 2025\\Issue: 25 February 2025};
\node[pair,right=of a] (b) {\textbf{Standard Chartered PLC, 7.625\%}\\Offering: 8 January 2025\\Issue: 16 January 2025};
\node[pair,below=of a] (c) {Retained public US\\distribution terms\\Minimum: USD 200,000};
\node[pair,below=of b] (d) {Retained US QIB or\\non-US offshore terms\\Minimum: USD 200,000};
\node[pair,below=of c] (e) {Apply the selected public-offer\\purchaser condition to\\the declared individual};
\node[pair,below=of d] (f) {US QIB condition: fails\\Non-US offshore condition: fails\\Same individual and account};
\node[pair,below=of e] (g) {\textbf{Selected condition: PASS}\\No failure of this purchaser\\condition in the declared model};
\node[pair,below=of f] (h) {\textbf{Selected condition: FAIL}\\Do not use this original\\distribution route for this buyer};
\node[mapwide,text width=132mm,below=8mm of $(g.south)!0.5!(h.south)$] (i) {\hyperref[map:coco-assessment]{\textbf{Carry both results into the complete bank assessment}}\\A purchaser condition does not decide the remaining layers};
\foreach \x/\y in {a/c,b/d,c/e,d/f,e/g,f/h}\draw[maparr](\x)--(\y);
\draw[maparr](g.south)--++(0,-4mm)-|(i.north);
\draw[maparr](h.south)--++(0,-4mm)-|(i.north);
\end{tikzpicture}
\end{MapDiagram}
\MapCaption{Both are real original offerings with the same minimum denomination.
The difference is the selected purchaser restriction. Their offer dates differ;
the comparison does not assert simultaneous live availability.}

Both instruments can expose holders to conversion and regulatory loss
absorption. A favourable purchaser result therefore says nothing about
principal protection. The two prospectuses and the retained 2025 legal
definitions support explicit interpretation proposals; page and file hashes
identify that evidence without proving the proposals correct.

Changing the actual purchaser can change the answer. A QIB buying for itself
can satisfy Standard Chartered's selected condition. A QIB bank acting for a
non-QIB client cannot use its own status to satisfy that client-account
condition. Moving an order offshore does not by itself remove US-person or
US-benefit restrictions. Section~\ref{sec:coco-purchaser-cases} derives the
formal rule and identifies the exact instruments and source passages.

\clearpage
\MapHeading{The CoCo example: the complete bank result}{map:coco-assessment}
The conditional purchaser finding supplements the existing bank investigation;
it does not close the broader product-law obligation. A \emph{qualified} result
explains why the evidence cannot support a complete transaction decision.

\begin{MapDiagram}
\begin{tikzpicture}[node distance=5mm and 8mm]
\tikzset{pair/.style={mapbox,text width=62mm,minimum height=19mm}}
\node[pair] (a) {Barclays\\Selected purchaser condition passes};
\node[pair,right=of a] (b) {Standard Chartered\\Selected purchaser condition fails};
\node[mapwide,text width=132mm,below=6mm of $(a.south)!0.5!(b.south)$] (c) {\textbf{Preserve all seven categories and fourteen obligations}\\Scope, sanctions, bank, client, product, internal policy, operations};
\node[mapwide,text width=132mm,below=of c] (d) {\textbf{Keep the actual missing evidence visible}\\Booking entity and service authority; client and beneficial owners;\\sanctions identity; other product rules; policy and settlement route};
\node[pair,below=6mm of d,xshift=-38mm] (e) {\textbf{Barclays: QUALIFIED}\\Purchaser pass retained;\\remaining gaps prevent clearance};
\node[pair,below=6mm of d,xshift=38mm] (f) {\textbf{Standard Chartered: QUALIFIED}\\Same unresolved bank layers;\\additional declared-route failure};
\node[mapwide,text width=132mm,below=6mm of $(e.south)!0.5!(f.south)$] (g) {\textbf{Neither result authorizes transaction execution}\\Preserve reasons and sources; reassess on changed evidence};
\draw[maparr](a.south)--++(0,-3mm)-|(c.north);
\draw[maparr](b.south)--++(0,-3mm)-|(c.north);
\draw[maparr](c)--(d);
\draw[maparr](d.south)--++(0,-3mm)-|(e.north);
\draw[maparr](d.south)--++(0,-3mm)-|(f.north);
\draw[maparr](e.south)--++(0,-3mm)-|(g.north);
\draw[maparr](f.south)--++(0,-3mm)-|(g.north);
\end{tikzpicture}
\end{MapDiagram}
\MapCaption{The implemented attachment preserves the bank investigation and adds
the purchaser finding. It cannot turn a partial pass into bank or regulatory
approval. The negative result concerns the declared original distribution.}

Issuer law is part of the product investigation even when the prospectus
contains no decisive contractual trigger. The CoCo survey in
Chapter~\ref{ch:circular} follows Swiss regulatory write-down powers and their
dated legal and litigation context. That survey informs required source checks;
it is not a completed executable model of every issuer jurisdiction. Swiss
provisions cannot be applied to these UK issues merely because all are CoCos.

Sanctions require their own analysis of the issuer, beneficial ownership,
program, activity and date. Chinese state ownership alone is not the selected
sanctions trigger. A name match supplies a candidate identity; no match does
not supply clearance. Restrictions on a US bank acting for a client also need
the bank's actual capacity and the client's status. These questions remain
separate from the US offering categories in the preceding chart.

The paired run reuses retained interpretation proposals; it does not conduct a
fresh model investigation of every prospectus clause. The
\hyperref[map:certificates]{formal and execution checks} support the stated
calculation without proving its English interpretation or supplying missing
bank evidence.

\clearpage
\MapHeading{Join the product conditions to the bank investigation}{map:joined-product}
The joined calculation keeps contractual loss, issuer powers and selling
requirements distinct. Section~\ref{sec:joined-eligibility} explains the
implemented conditions and their remaining premises.

\begin{MapDiagram}
\begin{tikzpicture}[node distance=5mm]
\node[mapwide] (a) {Retained prospectus and dated legal sources\\Identify the issue, mechanism, service and purchaser};
\node[mapbox,text width=56mm,below=7mm of a,xshift=-34mm] (b) {Contract and issuer powers\\Write-down or conversion;\\authority, order, implementation};
\node[mapbox,text width=56mm,below=7mm of a,xshift=34mm] (c) {Hong Kong product scope\\Direct debt or related product;\\legal form and exclusions};
\node[mapwide,below=7mm of $(b.south)!0.5!(c.south)$] (d) {Specific intermediary arrangement\\Check every broker-exception condition; retain missing facts};
\node[mapwide,below=of d] (e) {PI restriction and separate duties\\Exemptions do not erase the PI test; portfolio handling does not prove performance};
\node[mapwide,below=of e] (f) {Recompute the complete declared bank inventory\\All seven categories and fourteen obligations remain present};
\node[mapwide,below=of f] (g) {Qualified result with exact dependencies\\No execution permission from a partial pass; changes require reassessment};
\draw[maparr](a.south)--++(0,-3mm)-|(b.north);
\draw[maparr](a.south)--++(0,-3mm)-|(c.north);
\draw[maparr](b.south)--++(0,-3mm)-|(d.north);
\draw[maparr](c.south)--++(0,-3mm)-|(d.north);
\draw[maparr](d)--(e);\draw[maparr](e)--(f);\draw[maparr](f)--(g);
\end{tikzpicture}
\end{MapDiagram}
\MapCaption{A conversion can cancel debt principal while delivering shares.
Neither zero debt principal nor a PI classification settles the entire bank decision.}

The retained UBS issue converts under its specified conditions. The next chart
follows its implemented notice and share calculations. Swiss regulatory provisions do not
establish powers over the UK issues. The actual account, mandate, bank entities
and transaction route also remain necessary evidence. The joined reports for
all three prospectuses retain these gaps.

\clearpage
\MapHeading{Follow one UBS issue from publication to shares}{map:ubs-terms}
Section~\ref{sec:ubs-terms} uses Annex A of the June 2024 SGD notes. Source
proposals populate the product analysis; actual client and event facts remain
separate inputs.

\begin{MapDiagram}
\begin{tikzpicture}[node distance=5mm]
\node[mapwide] (a) {Identify the exact issue and source edition\\Retain terms, incorporated dependencies and amendments};
\node[mapbox,text width=56mm,below=7mm of a,xshift=-34mm] (b) {Capital-ratio branch\\Add higher-trigger capital;\\compare strictly below 7 percent};
\node[mapbox,text width=56mm,below=7mm of a,xshift=34mm] (c) {Viability branch\\FINMA and support premises;\\check alternative-loss notice};
\node[mapbox,text width=56mm,below=of b] (d) {Ordinary or extraordinary notice\\Restoration exception;\\higher-trigger postponement};
\node[mapbox,text width=56mm,below=of c] (e) {Three calendar days for notice\\Twenty Business Days\\from the viability event};
\node[mapwide,below=7mm of $(d.south)!0.5!(e.south)$] (f) {Retrieve actual notices and covered calendars\\A deadline does not establish an event};
\node[mapwide,below=of f] (g) {Reconstruct the price through the notice date\\Qualify missing history, source anomalies and unresolved rounding};
\node[mapwide,below=of g] (h) {Aggregate one holder's principal, then discard share fractions\\Check share creation and depository receipt separately};
\node[mapwide,below=of h] (i) {Join HKMA conditions and all fourteen bank obligations\\Preserve missing client, legal, sanctions and transaction evidence};
\draw[maparr](a.south)--++(0,-3mm)-|(b.north);
\draw[maparr](a.south)--++(0,-3mm)-|(c.north);
\draw[maparr](b)--(d);\draw[maparr](c)--(e);
\draw[maparr](d.south)--++(0,-3mm)-|(f.north);
\draw[maparr](e.south)--++(0,-3mm)-|(f.north);
\draw[maparr](f)--(g);\draw[maparr](g)--(h);\draw[maparr](h)--(i);
\end{tikzpicture}
\end{MapDiagram}
\MapCaption{The printed extraordinary-distribution formula has an unresolved
denominator anomaly. Its literal calculation is retained; no replacement price
or bank clearance is inferred.}

\clearpage
\MapHeading{Compare conversion notes with preferred shares}{map:preferred-evidence}
Section~\ref{sec:preferred-evidence} follows the Series SS preferred share.
Its non-cumulative dividend and redemption terms require different calculations
from the UBS conversion note. Both still enter the same bank assessment.

\begin{MapDiagram}
\begin{tikzpicture}[node distance=6mm]
\node[mapwide] (a) {Identify the issue, controlling terms and missing agreements\\Check source content, dates and amendment coverage};
\node[mapbox,text width=56mm,below=8mm of a,xshift=-34mm] (b) {UBS conversion note\\Event, notice and price history\\Conditional share entitlement};
\node[mapbox,text width=56mm,below=8mm of a,xshift=34mm] (c) {Series SS preferred share\\Declared dividends\\Redemption or liquidation};
\node[mapbox,text width=56mm,below=of b] (d) {Separate exchange calendars\\Scoped price determinations\\Custody and registration};
\node[mapbox,text width=56mm,below=of c] (e) {Cash allocation and taxes\\Controlling certificate's dates\\Actual payment evidence};
\node[mapwide,below=8mm of $(d.south)!0.5!(e.south)$] (f) {Match instrument, client and bank\\Check the same service and route\\Reject stale facts; preserve conflicts};
\node[mapwide,below=of f] (g) {Retain all fourteen bank obligations\\Missing law, sanctions or account evidence\\keeps the result qualified};
\node[mapwide,below=of g] (h) {Changed source or transaction\\Reassess the affected conclusions\\Retrieval alone does not prove legal force};
\draw[maparr](a.south)--++(0,-3mm)-|(b.north);\draw[maparr](a.south)--++(0,-3mm)-|(c.north);
\draw[maparr](b)--(d);\draw[maparr](c)--(e);
\draw[maparr](d.south)--++(0,-3mm)-|(f.north);\draw[maparr](e.south)--++(0,-3mm)-|(f.north);
\draw[maparr](f)--(g);\draw[maparr](g)--(h);
\end{tikzpicture}
\end{MapDiagram}
\MapCaption{Different instruments use different contractual calculations.
Evidence about a distribution, conversion or delivery remains separate from
permission for the bank to carry out the transaction.}

\clearpage
\MapHeading{1. Retain the source documents}{map:source}
The first task is to give every later reading the same inspectable source.
For PDFs, \texttt{pypdf} and Poppler's \texttt{pdftotext} are two different text
extractors. Web pages have two parsing paths. Comparing their text can expose
a missing line; checking the page for images can expose material that neither
extractor has read.

\begin{MapDiagram}
\begin{tikzpicture}[node distance=7mm and 5mm]
\node[mapbox] (a) {Source document\\Circular, law or prospectus};
\node[mapbox,right=of a] (b) {Keep the original\\Locate pages and passages};
\node[mapbox,right=of b] (c) {Extract the text\\Two extraction paths};
\node[mapbox,below=of b] (d) {Compare the extractions\\Inspect visual warnings};
\node[mapbox,below=of a] (e) {Retain usable text\\with its source locations};
\node[mapbox,below=of c] (f) {Unread or inconsistent?\\Keep the gap visible};
\draw[maparr] (a)--(b);\draw[maparr] (b)--(c);
\draw[maparr] (c.south)--++(0,-3mm)-|(d.north);
\draw[maparr] (d)--(e);
\draw[maparr] (d)--(f);
\end{tikzpicture}
\end{MapDiagram}
\MapCaption{Agreement between text extractors cannot reveal an exception present
only in an image. The additional source checks in Chapter~\ref{ch:ensemble}
read page images and compare the recognised text with the embedded text.}

\MapHeading{2. Fix context and applicable questions}{map:context}
A gift restriction may refer to the Code of Conduct and a frequently asked
question (FAQ). A retained authority catalog connects those references to
document editions and passages. A dependency search follows explicit official
links; missing documents remain named questions. The relevant business scope
identifies the people, products, activities and dates being assessed.

\begin{MapDiagram}
\begin{tikzpicture}[node distance=7mm and 5mm]
\node[mapbox] (a) {Circular passage\\and cited references};
\node[mapbox,right=of a] (b) {Authority catalog\\Edition and passage lookup};
\node[mapbox,right=of b] (c) {Date and scope check\\What can this source support?};
\node[mapbox,below=of b] (d) {Source set for the review\\with missing or uncertain items};
\draw[maparr] (a)--(b);\draw[maparr] (b)--(c);
\draw[maparr] (c.south)--++(0,-3mm)-|(d.north);
\end{tikzpicture}
\end{MapDiagram}
\MapCaption{Finding the FAQ changes the question from where the source is to
what its dated answer supports.}

The retained FAQ page has later answer updates. Its older page date therefore
does not establish everything a historical reader could have seen. Prospectus
work adds issue-specific supplements, incorporated documents, issuer law,
regulatory notices and the applicable selling rules. The transaction registry
binds factual documents, policy and settlement-route records to their retained
bytes and declared dates. Acquisition time does not establish legal force.
Missing or conflicting evidence remains explicit. Chapter~\ref{ch:circular}
develops these source questions; Appendix~\ref{app:current-interpretation} gives
the source tools.

\clearpage
\MapHeading{3. Develop possible interpretations}{map:readings}
A language model can propose what a passage requires and which facts would
decide a case. Separate requests produce initial readings without showing each
request the other answers. Another set of requests inventories the source's
requirements and qualifications. This second view helps expose an exception
that all the proposed readings might miss.

\begin{MapDiagram}
\begin{tikzpicture}[node distance=7mm and 5mm]
\node[mapbox] (a) {Source passages\\Review question};
\node[mapbox,right=of a] (b) {Separate model requests\\Readings and facts};
\node[mapbox,right=of b] (c) {Alternative readings\\with source reasons};
\node[mapbox,below=of b] (d) {Source inventory\\Duties and qualifications};
\node[mapbox,below=of c] (e) {Questions to investigate\\Differences and possible gaps};
\draw[maparr] (a)--(b);\draw[maparr] (b)--(c);
\draw[maparr] (a.south)--++(0,-3mm)-|(d.north);
\draw[maparr] (d)--(e);\draw[maparr] (c)--(e);
\end{tikzpicture}
\end{MapDiagram}
\MapCaption{Several readings reveal disagreement. A separate source inventory
can reveal a shared omission. Separate requests can still share mistakes.}

\MapHeading{4. Challenge the interpretations}{map:challenge}
Tree search organises successive revisions of a reading. Breadth-first search
visits the shallower alternatives first. Upper Confidence bounds applied to
Trees (UCT) balances previously useful branches against less-visited ones.
These are ways to choose the next investigation, not probabilities of legal
correctness. Both are implemented with finite work limits.

\begin{MapDiagram}
\begin{tikzpicture}[node distance=8mm and 5mm]
\node[mapbox] (a) {Select a reading\\Breadth-first or UCT};
\node[mapbox,right=of a] (b) {Challenge a reading\\Sources, examples, cases};
\node[mapbox,right=of b] (c) {Record the consequence\\Revision, rejection or open question};
\node[mapbox,below=of b] (d) {Keep alternatives and evidence\\Stop at the work limit};
\draw[maparr] (a)--(b);\draw[maparr] (b)--(c);
\draw[maparr] (c.north)--++(0,5mm)-|node[pos=.25,above,mapnote]{further work justified}(a.north);
\draw[maparr] (c.south)--++(0,-3mm)-|(d.north);
\end{tikzpicture}
\end{MapDiagram}
\MapCaption{Search changes a proposed reading only with an explicit reason.
Exhausting the work allowance leaves an unresolved result.}

Typed arguments record support and objections; official examples retain their
stated meaning. A rebate-and-voucher case can reveal where readings attach an
exception. After aligning facts and outputs, the original programs execute
that case. Draft compilation and step~6 therefore occur inside this loop too.
Agreement still cannot establish complete source coverage.
Chapters~\ref{ch:search} and~\ref{ch:ensemble} develop these checks.

\clearpage
\MapSubheading{4 continued. Investigate the exact disagreement}
The continuing investigation separates source support, relevance to the selected
question, correspondence with the executable rule and availability of required
authority. Two initial contexts work without peer answers. A later round can
then address the precise disagreement, or repair an invalid response, within
the same fixed allowance.

\begin{MapDiagram}
\begin{tikzpicture}[node distance=6mm]
\node[mapwide] (q) {Fix the question\\Actor, assessment unit, time and meaning of the result};
\node[mapwide,below=of q] (p) {Two blind initial assessments\\Source support, relevance, code correspondence, authority};
\node[mapwide,below=of p] (v) {Check identities, quotations and required evidence\\Retain both proposed explanations};
\node[mapwide,below=of v] (r) {Disagreement or missing evidence?\\Investigate the smallest consequential question};
\node[mapwide,below=of r] (o) {Retain supported proposals and unresolved alternatives\\Declared round limit; every dispatch consumes allowance};
\foreach \x/\y in {q/p,p/v,v/r,r/o}\draw[maparr](\x)--(\y);
\draw[maparr](r.east)--++(12mm,0)|-node[pos=.25,mapnote,rotate=90]{bounded follow-up}(p.east);
\end{tikzpicture}
\end{MapDiagram}
\MapCaption{A repaired response is another proposal. Agreement cannot supply a
missing source or prove an interpretation.}

Rival factual models retain different actors, objects and assessment times.
A proposed mapping permits a conditional comparison only after those
distinctions and its assumptions are stated. Changed source bytes or methods
start a new revision; the old program remains evidence for its old source.
Section~\ref{sec:successor-assurance} follows the implemented distinctions
through a contact history for a suspicious transaction report and the
unfamiliar-circular study.

Capacity is checked separately from meaning. Smaller source or comparison
batches retain the complete set of obligations and the original retry counts.
The optional incremental check admits each validated part only while total
retained size remains within its limit. A failed repair preserves earlier
claims; it cannot replace them with an empty inventory. If the work still
cannot finish, its missing checks remain in the final report.

\clearpage
\MapHeading{5. Construct one shared model before choosing a target}{map:rules}
A program needs the meaning of each fact and operation fixed. The shared
frontend records sources, proposed interpretations, typed facts, scoped rules
and uncertainty before a target language is selected. Type and dependency
checks expose invalid combinations and missing or circular references.
Separate translators then report which parts they can express.

\begin{MapDiagram}
\begin{tikzpicture}[node distance=8mm]
\node[mapwide,text width=132mm] (a) {\textbf{Shared frontend}\\Source evidence, proposed meanings, factual observations and scope};
\node[mapwide,text width=132mm,below=of a] (b) {\textbf{Shared typed model}\\Rules, dependencies, assessment times and uncertainty policy};
\node[mapbox,text width=60mm,below=of b,xshift=-38mm] (c) {\textbf{RuleIR translator}\\Check the declared capability profile};
\node[mapbox,text width=60mm,below=of b,xshift=38mm] (d) {\textbf{Native Catala translator}\\Check the declared capability profile};
\node[mapbox,text width=60mm,below=of c] (e) {RuleIR rules\\Java emitter and runtime};
\node[mapbox,text width=60mm,below=of d] (f) {Catala program\\Catala compiler and execution};
\node[mapwide,text width=132mm,below=8mm of $(e.south)!0.5!(f.south)$] (g) {\hyperref[map:java]{\textbf{Execute supported models with the same declared facts}}\\Unsupported translation stays an explicit capability failure};
\draw[maparr](a)--(b);
\draw[maparr](b.south)--++(0,-4mm)-|(c.north);
\draw[maparr](b.south)--++(0,-4mm)-|(d.north);
\draw[maparr](c)--(e);\draw[maparr](d)--(f);
\draw[maparr](e.south)--++(0,-4mm)-|(g.north);
\draw[maparr](f.south)--++(0,-4mm)-|(g.north);
\end{tikzpicture}
\end{MapDiagram}
\MapCaption{RuleIR is one target of the shared model. Native Catala does not
need to pass through RuleIR. A different proposed meaning requires a different
model, not merely a different target.}

For example, a gift restriction and a separately specified disclosure duty can
both belong in a package. Two rival readings of the restriction remain
alternative packages. The proposed grouping and each unresolved question are
preserved. There are two Catala routes. The compatibility route
lowers a supported RuleIR profile and retains the existing result, trace and
release checks. The native route consumes the shared model directly, so rich
records, lists, options, payload variants and typed operations can be reported
as executable Catala even when the RuleIR target returns an explicit capability
failure. These routes share the frontend and evidence policy; they do not make
different interpretations agree by construction.

\clearpage
\MapSubheading{5 continued. Where Stipula belongs}
Stipula describes contracts through parties, resources, actions and state
changes. It is relevant when consent can be granted and withdrawn. The
published translation studied in Chapter~\ref{ch:languages} constructs a Java
verification model with specifications in the Java Modeling Language (JML).
KeY is a verifier that reasons about the Java model and those specifications.

\begin{MapDiagram}
\begin{tikzpicture}[node distance=5mm]
\node[mapresearch] (a) {Stipula contract\\Parties and state changes};
\node[mapresearch,right=of a] (b) {Java model with JML\\Conditions and invariants};
\node[mapresearch,right=of b] (c) {KeY verification\\Named properties of that model};
\draw[maparr,dashed](a)--(b);\draw[maparr,dashed](b)--(c);
\end{tikzpicture}
\end{MapDiagram}
\MapCaption{Dashed boxes show a studied research route, not the current
LegalMath compiler. Its Java model is not automatically a bank application.}

The retained translator examples contain a compilation discrepancy; no local
KeY replay establishes their reported proofs. Adoption requires reproducible
inputs and a justified mapping to LegalMath's event rules.
Chapter~\ref{ch:languages} details those limits; Chapter~\ref{ch:verification}
explains the implemented rule language.

\clearpage
\MapHeading{6. Generate executable targets and check their behaviour}{map:java}
The Java emitter embeds the fixed RuleIR rules in a generated policy class
that calls a shared Java evaluator. The Java compiler, \texttt{javac}, compiles
that class and its supporting runtime. A separately written Python evaluator
supplies another execution of the same proposed rule. The Catala translators
take the same shared model when their capability report permits it: the
compatibility route lowers supported RuleIR, while the native route emits richer
Catala directly. Each target receives the same declared facts and assessment
times for the cases in its profile.

\begin{MapDiagram}
\begin{tikzpicture}[node distance=7mm]
\node[mapwide,text width=132mm] (a) {\textbf{Same model, facts, evidence identities and assessment times}\\Verify retained builds against the selected shared model};
\node[mapbox,text width=60mm,below=of a,xshift=-38mm] (b) {Execute RuleIR / Java\\Declared input profile};
\node[mapbox,text width=60mm,below=of a,xshift=38mm] (c) {Execute native Catala\\Same declared input profile};
\node[mapwide,text width=132mm,below=7mm of $(b.south)!0.5!(c.south)$] (d) {\textbf{Compare with separately written formal calculations}\\Check values and required abstentions; retain evidence and traces};
\node[mapbox,text width=60mm,below=of d,xshift=-38mm] (e) {Mismatch or invalid evidence\\Retain the failure; repair\\the affected model or implementation};
\node[mapbox,text width=60mm,below=of d,xshift=38mm] (f) {Agreement in the declared domain\\Retain the conditional result\\and remaining legal questions};
\draw[maparr](a.south)--++(0,-3.5mm)-|(b.north);
\draw[maparr](a.south)--++(0,-3.5mm)-|(c.north);
\draw[maparr](b.south)--++(0,-3.5mm)-|(d.north);
\draw[maparr](c.south)--++(0,-3.5mm)-|(d.north);
\draw[maparr](d.south)--++(0,-3.5mm)-|(e.north);
\draw[maparr](d.south)--++(0,-3.5mm)-|(f.north);
\end{tikzpicture}
\end{MapDiagram}
\MapCaption{The Java/Python pair and the Catala routes check whether a proposed
rule survives translation within a declared profile. They can agree on a rule
that omits a legal exception.}

The bank and CoCo runs use a complete-input profile: unknown, conflicting or
expired required evidence causes abstention. In the paired run, each target
returned 138 complete-input results and 19 abstentions. A successful abstention
test demonstrates preserved uncertainty, not a favourable purchaser finding.

The test cases must also be capable of exposing mistakes. Boundary cases place
values on either side of a threshold. Missing-fact cases check that uncertainty
survives. Event histories exercise consent, withdrawal and timing. Deliberately
altered rules, called \emph{mutations}, ask whether the tests notice a removed
exception or a reversed condition.

For a small finite rule, every declared input combination can be enumerated.
Larger or unsupported domains require more limited conclusions. Expected legal
outcomes need their own source justification, separate from Python/Java
agreement. Compilation failures and disagreements return to the affected rule
or implementation; successful checks accompany the draft rather than remove
its open interpretation questions. Chapter~\ref{ch:verification} develops
these checks and the separate Java host demonstration.

\clearpage
\MapHeading{6 continued. Where mathematical verification fits}{map:formal}
A mathematical check needs a precise proposition: whether two encoded rules
differ on a permitted input, for example, or whether a transfer preserves a
quantity. Such propositions concern the formal model and its assumptions.

\MapSubheading{Symbolic comparison with Z3}
Z3 is a satisfiability modulo theories (SMT) solver: it searches for values that
make a collection of logical and arithmetic conditions true. Here the condition
asks whether two supported rule encodings return different answers. The facts,
output meanings and permitted input domain must first be made comparable.

\begin{MapDiagram}
\begin{tikzpicture}[node distance=8mm and 5mm]
\node[mapbox] (a) {Comparable rules\\Supported input domain};
\node[mapbox,right=of a] (b) {Encode a disagreement\\Ask Z3 for a case};
\node[mapbox,right=of b] (c) {Different answers found\\Replay in Python and Java};
\node[mapbox,below=of a] (d) {Unsupported, unknown\\or inconsistent domain};
\node[mapbox,below=of b] (e) {No counterexample\\Within the stated domain};
\draw[maparr](a)--(b);\draw[maparr](b)--(c);
\draw[maparr](b)--node[right,mapnote]{ruled out}(e);
\draw[maparr](a)--(d);
\draw[maparr](b.south west)--(d.north east);
\end{tikzpicture}
\end{MapDiagram}
\MapCaption{A solver result concerns an encoded comparison. An unsupported
operation, timeout or empty domain cannot count as agreement.}

The solver adapter covers a stated subset of the formal model and does not
check an independent proof certificate from the solver. A found discrepancy is
valuable because the original programs can replay it. A no-counterexample
result retains the encoding and domain assumptions. Unsupported operations,
timeouts and empty domains remain explicit outcomes. Chapter~\ref{ch:proof}
explains those restrictions.

\clearpage
\MapHeading{6 continued. Prove a translation and challenge the tests}{map:certificates}
\MapSubheading{Proofs of selected mathematical properties}
Lean is a proof assistant that checks a formal derivation. The mathematical
review used it for selected logical and arithmetic propositions, including
three-valued logic, an exception counterexample and conservation under a stated
transfer. SymPy, a symbolic algebra system, checked selected algebraic identities.

\begin{MapDiagram}
\begin{tikzpicture}[node distance=5mm]
\node[mapbox] (a) {State a property\\Definitions and assumptions};
\node[mapbox,right=of a] (b) {Lean checks the proof\\Selected propositions};
\node[mapbox,right=of b] (c) {Retain the conclusion\\Only the property proved};
\draw[maparr](a)--(b);\draw[maparr](b)--(c);
\end{tikzpicture}
\end{MapDiagram}
\MapCaption{These are checks of selected mathematical claims, not an automatic
proof of each generated program or of the whole compiler.}

The successor also registers a narrow Lean certificate for Boolean expression
lowering. Its fixed checker proves the relation between two formal descriptions;
an additional finite check executes the generated Java for every declared
Boolean state assignment within its size limit. Arithmetic, dates, full event
histories and English meaning remain outside that theorem. A separate Java
entry point demonstrates attributed contact histories, keeping the duty trigger
distinct from an observed contact and from overall compliance.

For the CoCo purchaser rule, a Lean certificate checks preservation from the
formal expression to the shared model. The independent solver comparison covers
all 512 assignments of the nine Boolean premises. Native execution is a
separate check: both targets receive the same 157 complete, uncertain or
expired cases. These are different forms of evidence about a stated model.

\MapSubheading{Would the tests detect a consequential mistake?}
The purchaser tests deliberately remove conditions that matter. One mutation
ignores whose account benefits; another ignores the purchaser's US-person
status. A revealing input must make the altered rule disagree with the
independently expressed relation.

\begin{MapDiagram}
\begin{tikzpicture}[node distance=5mm]
\node[mapbox] (a) {Introduce a defect\\Ignore the client account};
\node[mapbox,right=of a] (b) {Find a revealing case\\QIB bank; non-QIB client};
\node[mapbox,right=of b] (c) {Retain the counterexample\\Did the test expose the defect?};
\draw[maparr](a)--(b);\draw[maparr](b)--(c);
\end{tikzpicture}
\end{MapDiagram}
\MapCaption{Six altered purchaser rules produced counterexamples. Mutation
testing checks sensitivity to those defects; it does not manufacture legal
answers or establish that the source interpretation is complete.}

Proving that every supported rule preserves its behaviour through translation
remains a stronger goal. The present checks do not establish that whole-compiler
result or correct interpretation of English. The Stipula/JML/KeY route offers
another research approach to particular Java properties.

\clearpage
\MapHeading{The proof-carrying dossier and the abstention path}{map:assurance}
The preceding checks are useful only when their targets and assumptions stay
visible. The integration record therefore binds a source packet, competing
hypotheses, arguments, proof obligations, method-family results and executable
artifacts by hash. It also records shared dependencies: agreement between Java
and Catala after the same RuleIR lowering is a translation check, not two
independent readings. The controller never changes an unavailable, unknown or
unsupported result into agreement.

\begin{MapDiagram}
\begin{tikzpicture}[node distance=18mm and 7mm]
\tikzset{mapcompact/.style={mapbox,text width=37mm,minimum height=30mm,inner sep=2.5mm,font=\small}}
\node[mapcompact] (a) {Versioned source packet\\Question and scope\\Quoted passages};
\node[mapcompact,right=of a] (b) {Interpretation families\\Arguments and search\\Open assumptions};
\node[mapcompact,right=of b] (c) {Formal obligations\\Solver/proof checks\\Java and Catala artifacts};
\node[mapcompact,below=of b] (d) {Verify retained files\\Targets and hashes\\Common inputs};
\node[mapcompact,below=of a] (e) {Unresolved or novel\\clause\\Extension required};
\node[mapcompact,below=of c] (f) {Reviewable draft\\Evidence dossier\\No release authority};
\draw[maparr] (a)--(b);\draw[maparr] (b)--(c);\draw[maparr] (c)--(d);
\draw[maparr] (d)--(e);\draw[maparr] (d)--(f);
\draw[maparr] (e.north)--++(0,9mm)-|node[pos=.25,above,mapnote]{new evidence or revised model}(b.south);
\end{tikzpicture}
\end{MapDiagram}
\MapCaption{A proof or replay closes the proposition it names. An open authority,
future clause, stale result or unavailable required method keeps the dossier on
the abstention path.}

The gift-circular replay (23EC46) imports five previously proposed readings. Two differ
on whether product-directed presentation suffices without an operational
connection. The solver finds a separating assignment and the original Python
and Java programs replay it. The Catala routes supply another execution target
over the shared model: the compatibility backend consumes supported RuleIR,
while the native converter can consume a richer model without first making it
RuleIR. Agreement after a common lowering is a translation check, not
independent legal evidence. The disabled live provider supplies no fresh
reading. The witness identifies a legal question; it does not answer it.
The result can be delivered
as a draft Java package with its evidence and open questions. The transaction
assessment must retain the unresolved classification and any missing policy.
The exact contract and master command are described in
Section~\ref{sec:proof-carrying-integration}.

The master validates a frozen revision and binds each accepted phase to its
inputs and predecessor files.
Failure preserves the attempt and triggers a focused repair check before a
bounded retry. Success refreshes the next-phase plan. A changed source, checker,
plan or allowlist invalidates acceptance; retry exhaustion produces a report of
the remaining failure. Neither a retry nor a successful document build supplies
missing interpretive evidence.

\clearpage
\MapHeading{7. Assemble the transaction assessment}{map:delivery}
A Java archive (JAR) or a declared Catala package contains the compiled
control. Its execution result is only one part of the transaction assessment.
The investigation also checks the fixed obligation inventory, applicability,
retrieved factual documents and their source dependencies. Each finding keeps
its scope, supporting evidence and unresolved questions.

\begin{MapDiagram}
\begin{tikzpicture}[node distance=7mm]
\node[mapwide,text width=132mm] (a) {\textbf{Retrieve evidence for the declared transaction}\\Actual document fields; source versions, scope and time checks};
\node[mapwide,text width=132mm,below=of a] (b) {\textbf{Calculate supported conditions in dependency order}\\Derived facts cannot be replaced by an asserted favourable result};
\node[mapwide,text width=132mm,below=of b] (c) {\textbf{Reconcile every obligation in the fixed inventory}\\Findings, missing rules, conflicts and unsupported applicability};
\node[mapwide,text width=132mm,below=of c] (d) {\textbf{Bind the assessment to sources, facts and method}\\Program identity, formal checks, execution results and reasons};
\node[mapwide,text width=132mm,below=of d] (e) {\textbf{Current result: qualified assessment}\\Complete legal coverage is not established; execution is not authorized};
\foreach \x/\y in {a/b,b/c,c/d,d/e}\draw[maparr](\x)--(\y);
\end{tikzpicture}
\end{MapDiagram}
\MapCaption{A request supplies evidence and context, not its own shorter list
of requirements or precomputed favourable findings. The retained assessment
can be reconstructed from its bound inputs.}

For example, account amounts enter the financial condition before client
qualification is calculated. A transaction threshold also needs its separately
declared controls premise. Sophisticated professional investor eligibility then
feeds the condition for streamlined procedures. A false sufficient condition
does not prove its converse: failure to establish that a bond is complex does
not establish that it is non-complex.

The CoCo purchaser result is supplementary evidence attached to this assessment.
The original fourteen bank obligations remain present for each instrument.
Actual JPMorgan entity, account, policy and settlement facts have not been
supplied. The program therefore retains those gaps and returns no transaction
execution permission. Recording a bank authorization as a business fact would
not turn it into proof of legal interpretation quality.

\clearpage
\MapHeading{8. Reassess when evidence or events change}{map:operation}
An assessment belongs to particular source editions, factual records, a method
and an assessment time. It cannot be reused as though a later transaction had
the same circumstances. The event replay records changes and invalidates the
current assessment before it considers a later settlement request.

\begin{MapDiagram}
\begin{tikzpicture}[node distance=7mm]
\node[mapwide] (a) {\textbf{Retained assessment and ordered events}\\Exact sources, facts, policy, route and dates};
\node[mapwide,below=of a] (b) {\textbf{A relevant input changes}\\Law or sanctions source; account facts; policy; route; corporate action};
\node[mapwide,below=of b] (c) {\textbf{Invalidate the current assessment}\\Preserve its historical evidence; retrieve the changed inputs};
\node[mapwide,below=of c] (d) {\textbf{Reassess applicability and all declared obligations}\\Recheck affected models and their dependencies};
\node[mapwide,below=of d] (e) {\textbf{Retain the new result and any separately sourced duties}\\No current qualified result permits settlement};
\foreach \x/\y in {a/b,b/c,c/d,d/e}\draw[maparr](\x)--(\y);
\draw[maparr](e.east)--++(12mm,0)|-node[pos=.25,mapnote,rotate=90]{next change}(b.east);
\end{tikzpicture}
\end{MapDiagram}
\MapCaption{Settlement requires a current assessment for its declared time.
A stale or absent assessment is rejected. Passing that freshness check still
does not turn a qualified result into permission to trade.}

Duties keep the responsible actor, trigger, source and due time separate.
Rejection, blocking assets, reporting, rescreening and disclosure are different
actions; a sanctions finding cannot silently substitute one for another. A
business-day deadline needs an explicit retained calendar within its supported
range. Evidence recorded as completion does not prove the action occurred.

Official sanctions snapshots can be refreshed and checked before replacement;
earlier bytes remain available for historical investigations. Exact-name matches
identify candidates for entity resolution. Ownership propagation for a blocking
program cannot automatically be reused for a solely Chinese military-industrial
company designation. These scoped mechanisms do not complete identity resolution
or every sanctions program.

\clearpage
\MapHeading{Repair the implementation; refresh the next phase}{map:repair}
The executable master programs preserve the evidence from each phase. The bank
and paired-CoCo campaigns validate sources and tools, run regressions, check
formal relations, execute both targets, and then assemble the integrated
investigation. A phase failure must remain visible until a repair has been
executed and checked.

\begin{MapDiagram}
\begin{tikzpicture}[node distance=8mm]
\node[mapwide] (a) {\textbf{Freeze the phase inputs and required checks}\\Source, model, tool and predecessor identities};
\node[mapwide,below=of a] (b) {\textbf{Execute the current phase}\\Keep outputs and diagnostics under a new attempt};
\node[mapbox,text width=58mm,below=of b,xshift=-37mm] (c) {\textbf{Failure}\\Retain the failed attempt\\Record the required repair};
\node[mapbox,text width=58mm,below=of b,xshift=37mm] (d) {\textbf{Checks pass}\\Validate output identities\\Accept only this phase's result};
\node[mapbox,text width=58mm,below=of c] (e) {Execute and check the repair\\Rerun with current inputs\\Changed methods invalidate reuse};
\node[mapbox,text width=58mm,below=of d] (f) {Refresh the next-phase plan\\Continue to the next phase\\or retain the final qualified result};
\draw[maparr](a)--(b);
\draw[maparr](b.south)--++(0,-4mm)-|(c.north);
\draw[maparr](b.south)--++(0,-4mm)-|(d.north);
\draw[maparr](c)--(e);\draw[maparr](d)--(f);
\draw[maparr](e.west)--++(-6mm,0)|-(a.west);
\end{tikzpicture}
\end{MapDiagram}
\MapCaption{The runner records a repair request; a failed check does not repair
itself. Reuse of a successful phase requires matching inputs and output hashes.
An engineering pass does not remove the legal qualifications.}

\MapSubheading{Evaluate new situations without human answer labels}
The product requires a machine-checked proposition or an explicit qualification.
Human approval and agreement among models are not quality evidence. Independent
equations, translation proofs, adversarial cases and source-bound challenges
test different possible failures, each within its declared domain.

For future-source work, the observation mechanism freezes the method, inventory
and prior sources before it accepts new content. A changed method starts a new
series. New-to-inventory documents and generated amendments test novelty handling;
they do not establish accuracy on future laws. The present bank series has no
real prospective observations. Unknown meaning must remain qualified until
evidence establishes the particular claim being made.

\clearpage
\MapHeading{Complete the processing without losing the question}{map:continuation}
The same proof-or-qualification requirement applies when a source investigation
must be split into smaller pieces. Finishing those pieces matters because an
unexamined exception can invalidate a conclusion drawn from the rest.

Long responses and exhausted review budgets can leave parts of a circular
unexamined. Two implemented paths share exact source references. The single
reader accounts for its declared questions; the ensemble investigates competing
readings. These paths retain their separate results. See
Section~\ref{sec:processing-continuation} for the mechanics and checked limits.

\begin{MapDiagram}
\begin{tikzpicture}[node distance=5mm]
\node[mapwide] (a) {Whole retained source\\Select exact source identifiers; software retrieves the text};
\node[mapbox,text width=54mm,below=7mm of a,xshift=-34mm] (b) {Single reading\\Exact declared questions};
\node[mapbox,text width=54mm,below=7mm of a,xshift=34mm] (c) {Competing readings\\Exact question assignments};
\node[mapbox,text width=54mm,below=of b] (d) {Bounded complete coverage\\Every source unit accounted for};
\node[mapbox,text width=54mm,below=of c] (e) {Breadth before follow-ups\\Two perspectives per batch};
\node[mapbox,text width=54mm,below=of d] (f) {Whole-source consistency\\Recheck coverage\\after revision};
\node[mapbox,text width=54mm,below=of e] (g) {Scoped reconciliation\\Retain disputes and missing work};
\node[mapwide,below=7mm of $(f.south)!0.5!(g.south)$] (h) {Retained proposals and uncertainty\\Source, reading, question and method identities};
\draw[maparr](a.south)--++(0,-3mm)-|(b.north);
\draw[maparr](a.south)--++(0,-3mm)-|(c.north);
\draw[maparr](b)--(d);\draw[maparr](c)--(e);
\draw[maparr](d)--(f);\draw[maparr](e)--(g);
\draw[maparr](f.south)--++(0,-3mm)-|(h.north);
\draw[maparr](g.south)--++(0,-3mm)-|(h.north);
\end{tikzpicture}
\end{MapDiagram}
\MapCaption{An interpretation revision reopens its coverage checks. Failed
requests keep their spent attempts; no batch disappears to fit the budget.}

The qualification adapter checks supported formal propositions for a retained
reading and exact question. The complete ensemble investigation invokes it when
enabled; the split single-reader path retains proposals for subsequent checking.
A source identifier proves where the words came from. Coverage proves that the
declared units and questions received dispositions. A formal theorem proves its
stated relation. None of these alone proves the English interpretation. The
four earlier unfamiliar investigations therefore remain incomplete unless their
remaining stages actually execute under the recorded limits.

\clearpage
\MapHeading{Preserve every concern before shortening a request}{map:lossless-context}
The retained authentication exercise reached all 263 source units but could not
fit its final whole-source comparison into one request. The optional readable
table encoding removes repeated representation while retaining the complete
source, questions and explanations. Section~\ref{sec:lossless-context} explains
the encoding and its executed checks.

\begin{MapDiagram}
\begin{tikzpicture}[node distance=5mm]
\node[mapwide] (a) {Retain original complete request\\Full source, every question and all cross-piece concerns};
\node[mapwide,below=of a] (b) {Encode readable tables and exact references\\No summary, truncation or deleted qualification};
\node[mapwide,below=of b] (c) {Decode and compare the entire original request\\Canonical bytes equal; original identities unchanged};
\node[mapbox,text width=57mm,below=7mm of c,xshift=-34mm] (d) {Preserved and within byte limit\\Dispatch complete context};
\node[mapbox,text width=57mm,below=7mm of c,xshift=34mm] (e) {Mismatch or overflow\\Keep evidence; stop dispatch};
\node[mapbox,text width=57mm,below=of d] (f) {Whole-source reconciliation\\Account for every concern\\and question; retain uncertainty};
\node[mapbox,text width=57mm,below=of e] (g) {Bounded causal repair\\Execute its diagnostic\\before a permitted retry};
\node[mapwide,below=7mm of $(f.south)!0.5!(g.south)$] (h) {An interpretation revision reopens all coverage\\No discarded history or refunded attempts};
\draw[maparr](a)--(b);\draw[maparr](b)--(c);
\draw[maparr](c.south)--++(0,-3mm)-|(d.north);
\draw[maparr](c.south)--++(0,-3mm)-|(e.north);
\draw[maparr](d)--(f);\draw[maparr](e)--(g);
\draw[maparr](f.south)--++(0,-3mm)-|(h.north);
\draw[maparr](g.south)--++(0,-3mm)-|(h.north);
\end{tikzpicture}
\end{MapDiagram}
\MapCaption{Exact reconstruction permits the capacity comparison; it does not
establish legal support. The executed synthetic reconciliation retained every
concern as unresolved.}

The measured request shrank from 382,396 to 175,016 bytes under the unchanged
200,000-byte limit. A challenge with long unique explanations still exceeded
that limit and stopped before dispatch. A future source can therefore produce
an explicit capacity qualification. Model comprehension of the encoding and
English correctness remain unestablished by these offline exercises.

\clearpage
\MapHeading{Check program references and dates independently}{map:executable-dates}
Section~\ref{sec:executable-dates} separates exact quotation, mathematical
preservation and native execution. Each establishes a different proposition.

\begin{MapDiagram}
\begin{tikzpicture}[node distance=6mm]
\node[mapwide] (a) {Retained source, proposed reading and typed question\\Check the request against these original objects};
\node[mapbox,text width=56mm,below=7mm of a,xshift=-34mm] (b) {Parse scope and result\\Select their exact expressions\\Fact descriptions remain context};
\node[mapbox,text width=56mm,below=7mm of a,xshift=34mm] (c) {Reading is unencoded\\No executable quotation\\or correspondence verdict};
\node[mapwide,below=7mm of $(b.south)!0.5!(c.south)$] (d) {Check ownership; resolve text; retain all proposed judgments\\Run original semantic checks; preserve rejected batches};
\node[mapwide,below=of d] (e) {For supported Gregorian comparisons\\Check constructor preservation and calendar-order theorems};
\node[mapwide,below=of e] (f) {Check the actual representations and runtimes\\Independent calendar codec; RuleIR/Java and Catala challenges};
\node[mapwide,below=of f] (g) {Retain the legal date question and missing authorities\\Freeze the method before collecting future observations};
\draw[maparr](a.south)--++(0,-3mm)-|(b.north);
\draw[maparr](a.south)--++(0,-3mm)-|(c.north);
\draw[maparr](b.south)--++(0,-3mm)-|(d.north);
\draw[maparr](c.south)--++(0,-3mm)-|(d.north);
\draw[maparr](d)--(e);\draw[maparr](e)--(f);\draw[maparr](f)--(g);
\end{tikzpicture}
\end{MapDiagram}
\MapCaption{The calendar domain is valid proleptic Gregorian dates in years
1--9999. Business-day rules and selection of the legally relevant date remain separate.}

The eight encoded UCITS readings can now receive date-preservation evidence.
The ninth remains unencoded, and unresolved source questions persist for all
nine. Tests of this implementation do not create future observations or certify
the interpretation of an unknown future legal text.

Source extraction and question definition use source references. Version~2
executable references apply only at the correspondence stage. A failed response
triggers bounded diagnostic processing; exhaustion retains every unresolved
pair and concern. Exact quotation never establishes legal support.

\clearpage
\MapHeading{Keep the method fixed when observing future sources}{map:whole-method}
Section~\ref{sec:whole-method-observation} explains why a study must identify
the extraction and interpretation process as well as its final proof checker.

\begin{MapDiagram}
\begin{tikzpicture}[node distance=6mm]
\node[mapwide] (a) {Freeze the method before admitting new sources\\Extraction, model route, search and questions\\Authority, spending limits and proof tools};
\node[mapwide,below=of a] (b) {Record each submitted source edition and question\\Check publication and encounter dates; retain excluded and pending tasks};
\node[mapwide,below=of b] (c) {Execute the investigation with original spending limits\\Retain all candidates, rejected responses, discrepancies and missing premises};
\node[mapwide,below=of c] (d) {Join the actual stage records to the frozen method\\Recheck each candidate qualification and the complete pair denominator};
\node[mapbox,text width=56mm,below=7mm of d,xshift=-34mm] (e) {Method and records match\\Retain the checked propositions\\and every unresolved question};
\node[mapbox,text width=56mm,below=7mm of d,xshift=34mm] (f) {Method changed\\or stage missing\\Retain incomplete task\\Repair needs a later window};
\node[mapwide,below=8mm of f.south,xshift=-34mm] (g) {Report the observed evidence and its limits\\Legal meaning and remote attestation remain qualified\\Publication completeness is unestablished};
\draw[maparr](a)--(b);\draw[maparr](b)--(c);\draw[maparr](c)--(d);
\draw[maparr](d.south)--++(0,-3mm)-|(e.north);
\draw[maparr](d.south)--++(0,-3mm)-|(f.north);
\draw[maparr](e.south)--++(0,-3mm)-|(g.north);
\draw[maparr](f.south)--++(0,-3mm)-|(g.north);
\end{tikzpicture}
\end{MapDiagram}
\MapCaption{A checked program belongs to a particular investigation. The same
program produced from a changed source or method is a different observation.}

An additional allowance does not reset the original question's attempts or
deadline. A revised source keeps its predecessor's unresolved questions. A
repaired method needs a later untouched observation window; no earlier failure
is silently turned into evidence for that revision.

\clearpage
\MapHeading{Challenge a reading before eliminating it}{map:challenge-contract}
Section~\ref{sec:systematic-challenges} explains the distinction between an
unresolved source reading and a checked mathematical counterexample.

\begin{MapDiagram}
\begin{tikzpicture}[node distance=6mm]
\node[mapwide] (a) {Retain the full source, original question and every candidate};
\node[mapwide,below=of a] (b) {Generate declared challenges\\Actor, scope, quantifiers, negation, exceptions, time, definitions, authority and output};
\node[mapwide,below=of b] (c) {Execute different checking mechanisms\\Source cues, blind reader proposals and formal counterexamples};
\node[mapwide,below=of c] (d) {Join bounded answers with complete context\\Missing replies and substantive differences remain visible};
\node[mapbox,text width=56mm,below=7mm of d,xshift=-34mm] (e) {Replayed counterexample\\Candidate violates a stated\\formal obligation};
\node[mapbox,text width=56mm,below=7mm of d,xshift=34mm] (f) {Source meaning unresolved\\Retain rival readings and\\missing legal premises};
\node[mapwide,below=8mm of f.south,xshift=-34mm] (g) {Bounded follow-up or explicit uncertainty\\Preserve prior responses, candidates and spending};
\draw[maparr](a)--(b);\draw[maparr](b)--(c);\draw[maparr](c)--(d);
\draw[maparr](d.south)--++(0,-3mm)-|(e.north);
\draw[maparr](d.south)--++(0,-3mm)-|(f.north);
\draw[maparr](e.south)--++(0,-3mm)-|(g.north);
\draw[maparr](f.south)--++(0,-3mm)-|(g.north);
\end{tikzpicture}
\end{MapDiagram}
\MapCaption{A counterexample rejects a candidate relative to a declared premise.
It does not establish that the English source warrants that premise.}

The source scanner cannot prove that all exceptions were found. A finite grammar
cannot prove that every legal interpretation was considered. These qualifications
remain explicit when the checked stages are joined into one investigation.

\endgroup
\clearpage
